\thispagestyle{empty}
\vspace*{\fill}
\begin{center}
  \textbf{Notasi}  \\[2ex]
  {\renewcommand{\arraystretch}{1.15}%
  \begin{tabular}{@{}>{\raggedleft\arraybackslash}p{0.36\textwidth}|>{\raggedright\arraybackslash}p{0.56\textwidth}@{}}
    \( \Re \), \( \Re^+ \), \( \Re^n \) &bilangan real, bilangan real positif, tupel-$n$ bilangan real \\
    \( \N              \),
    \( \C              \)  &bilangan asli \( \set{0,1,2,\ldots} \), bilangan kompleks                           \\
    \( (a\,..\,b) \), \( [a\,..\,b] \) &interval terbuka, interval tertutup   \\
    \( \sequence{\ldots} \)&barisan (daftar yang urutannya diperhitungkan)    \\
    \( h_{i,j} \)          &entri baris \( i \) dan kolom \( j \) dari matriks~$H$ \\
    \( V,W,U \)            &ruang vektor               \\
    \( \vec{v} \),
    $\zero$, $\zero_V$     &vektor, vektor nol, vektor nol dari ruang~$V$   \\
    \( \polyspace_n \), \( \matspace_{\nbym{n}{m}} \)  
                          &ruang polinomial berderajat~\( n \), matriks \( \nbym{n}{m} \)      \\
    \( \spanof{S} \)       &rentang suatu himpunan                   \\
    \( \sequence{B,D} \), \( \vec{\beta},\vec{\delta} \)         
                          &basis, vektor basis  \\
    \( \stdbasis_n=\sequence{\vec{e}_1,\,\ldots,\,\vec{e}_n} \)         
                          &basis standar untuk $\Re^n$  \\
    \( V\isomorphicto W \) &ruang yang isomorfik                         \\
    \( M\directsum N \)    &jumlah langsung subruang                   \\
    \( h,g \)              &homomorfisme (pemetaan linear)                \\
    % \( H,G \)              &matrices                                  \\
    \( t,s \)              &transformasi (pemetaan linear dari suatu ruang ke dirinya sendiri) \\
    % \( T,S \)              &square matrices                           \\
    \( \rep{\vec{v}}{B} \), \( \rep{h}{B,D} \)     
                          &representasi suatu vektor, suatu pemetaan    \\
    \( Z_{\nbym{n}{m}} \) atau \( Z \), \(I_{\nbyn{n}}\) atau \( I \)  &matriks nol, matriks identitas    \\
    \( \deter{T} \)        &determinan matriks       \\
    \( \rangespace{h},\nullspace{h} \)
                           &ruang hasil, ruang nol dari pemetaan  \\
    \( \genrangespace{h},\gennullspace{h} \)
                           &ruang hasil tergeneralisasi dan ruang nol tergeneralisasi
  \end{tabular}
  }
\end{center}
\vspace*{\fill}
\clearpage
\thispagestyle{empty}
\vspace*{\fill}
\begin{center}
  \textbf{Huruf Yunani beserta pelafalannya}
    \\[1.5ex]
  \newcommand{\pronounced}[1]{\hspace*{.2em}\small\textit{#1}}
  {\renewcommand{\arraystretch}{1.25}%
  \begin{tabular}{cl@{\hspace*{3em}}cl}
    karakter &\multicolumn{1}{c}{\makebox[-3.5em][r]{nama}}       
    &karakter  &\multicolumn{1}{c}{\makebox[-3.5em][r]{nama}}  \\ 
    \hline
     \makebox[1em][l]{\( \alpha  \)} &alfa \pronounced{AL-fa}  
       &\makebox[1em][l]{\( \nu     \)}  &nu  \pronounced{NIU}       \\
     \makebox[1em][l]{\( \beta   \)} &beta  \pronounced{BE-ta}     
       &\makebox[1em][l]{\( \xi  \), \( \Xi \)}  &ksi   \pronounced{KSI}    \\ 
     \makebox[1em][l]{\( \gamma  \), \( \Gamma \)} &gama  \pronounced{GA-ma}
       &\makebox[1em][l]{\( o \)} &omikron  \pronounced{O-mi-kron}  \\
     \makebox[1em][l]{\( \delta  \), \( \Delta \)} &delta  \pronounced{DEL-ta} 
       &\makebox[1em][l]{\( \pi \), \( \Pi \)} &pi  \pronounced{PI}     \\
     \makebox[1em][l]{\( \epsilon\)} &epsilon  \pronounced{EP-si-lon}   
       &\makebox[1em][l]{\( \rho \)} &rho  \pronounced{RO}    \\
     \makebox[1em][l]{\( \zeta   \)} &zeta   \pronounced{ZE-ta}    
       &\makebox[1em][l]{\( \sigma  \), \( \Sigma \)} &sigma  \pronounced{SIG-ma}  \\
     \makebox[1em][l]{\( \eta  \)} &eta  \pronounced{E-ta}      
       &\makebox[1em][l]{\( \tau \)} &tau  \pronounced{TAU}    \\
     \makebox[1em][l]{\( \theta \), \( \Theta \)} &teta  \pronounced{TE-ta}    
       &\makebox[1em][l]{\( \upsilon\), \( \Upsilon \)} &upsilon  \pronounced{UP-si-lon}  \\
     \makebox[1em][l]{\( \iota \)} &iota \pronounced{I-o-ta}   
       &\makebox[1em][l]{\( \phi \), \( \Phi \)} &fi  \pronounced{FI}    \\
     \makebox[1em][l]{\( \kappa  \)} &kapa  \pronounced{KA-pa}  
       &\makebox[1em][l]{\( \chi \)}  &khi  \pronounced{KHI}    \\
     \makebox[1em][l]{\( \lambda \), \( \Lambda \)} &lambda  \pronounced{LAM-da}  
       &\makebox[1em][l]{\( \psi    \), \( \Psi \)}  &psi \pronounced{PSI}    \\
     \makebox[1em][l]{\( \mu  \)}  &mu  \pronounced{MIU}     
       &\makebox[1em][l]{\( \omega  \), \( \Omega \)} &omega  \pronounced{O-me-ga}  
  \end{tabular}}   \\[1.5ex]
  Huruf kapital yang ditampilkan adalah yang berbeda dari huruf kapital Romawi.
\end{center}
\vspace*{\fill}
